\section{讨论}

\subsection{为何“统一化”重要}

本文最重要的结果既是算法性的，也是概念性的。通过在目标响应场中同时纳入角度、频率与偏振三个维度，论文将傅里叶算子超表面设计重构为一个单一逆问题，而不同器件行为只是在这一统一对象上的不同投影。这一转变之所以重要，是因为它将讨论从若干孤立示例推进到对自由形态光学算子设计中“应该学习什么、生成什么、验证什么”的更一般性理解。

在这一意义上，物理一致生成式逆向设计并不是一个套在超表面问题上的通用机器学习外壳，而是一种与逆问题结构相匹配的建模选择。该问题本身具有多解性，目标是场而非标量，而物理验证始终不可省略。正是这些因素共同构成了本文框架的科学一致性。

\subsection{局限性与未来问题}

当然，当前表述仍有若干局限。首先，本文强调的是响应场匹配，尚未充分利用某些算子可能需要的更完整复场信息。其次，角度与波长的离散采样不可避免地限制了模型能够学习和验证的范围，尤其对于尖锐共振特征更是如此。第三，尽管自由形态二值参数化有利于展示统一逆向设计思想，但若面向特定工艺，仍需引入更强的制造约束。

这些局限更适合被视为在同一响应空间思想下进一步扩展的方向，而不是放弃这一框架的理由。未来工作可以在保持整体结构不变的前提下，引入相位敏感目标、不确定性感知条件，以及更强的工艺鲁棒先验。
